\section{The Fathers Do Not Speak with One Flat Voice}

\subsection{Ignatius and Irenaeus: incarnation against unreality}

Ignatius of Antioch does not furnish a commentary on John 6, but his
\emph{Letter to the Smyrnaeans} 7 identifies the Eucharist as the flesh of
Jesus Christ that suffered and was raised. The polemical target is docetism:
those who deny the reality of Christ's flesh abstain from Eucharist. This is
important reception evidence because it refuses to separate sacramental
confession from the flesh born, suffering, and raised. It does not establish
that Ignatius was expounding John 6 at that locus.

Irenaeus develops the same anti-docetic and anti-gnostic pressure. In
\emph{Against Heresies} IV.18.5, bread and wine from creation become the
Eucharist and nourish human flesh; in V.2.2--3, bodies nourished from Christ's
body and blood are destined for resurrection. The logic closely matches John
6's conjunction of flesh, eating, and resurrection, while remaining Irenaeus's
own argument against those who deny the Creator's salvation of embodied
creatures. Patristic convergence is not proof of a quotation.

\subsection{Cyprian and the Jerusalem mystagogical catechesis}

Cyprian's \emph{On the Lord's Prayer} 18 places John 6 inside an explicitly
Eucharistic account of daily bread. He names Christ as the bread of life,
describes daily reception as food of salvation, cites the flesh given for the
world and the command to eat and drink, and connects exclusion from communion
with separation from Christ's body.

The fourth Jerusalem mystagogical catechesis begins from the institution
narrative, cites John 6:53, rejects a carnal construal, and directs the newly
initiated to judge the Eucharistic gift by Christ's declaration rather than
taste. The traditional attribution is to Cyril of Jerusalem, but authorship of
the mystagogical lectures remains disputed; the checked witness establishes
Jerusalem catechesis and direct Eucharistic reception without resolving that
question.

\subsection{Chrysostom: faith, flesh, mystery, and union}

Chrysostom's Homilies 45--47 on John preserve the discourse's movement rather
than forcing one sense on every occurrence of ``bread.'' At 6:35, he relates
the bread of life first to Christ's divinity; at 6:51, he marks the transition
to the flesh Christ will give. He allows that ``bread'' can signify saving
teaching and faith or Christ's own body, because the context governs.

At the decisive turn, Chrysostom calls the manna type and shadow and Christ's
gift the reality. He reads the ``how'' of 6:52 as unbelief when it subjects
divine action to ordinary physical explanation. Yet his Eucharistic language
is not attenuated: Christ gives himself to be touched and eaten so that the
faithful become one body with him. In Homily 47, he reads 6:63 as a prohibition
on carnal hearing, explicitly denying that ``the flesh profits nothing'' means
Christ's own flesh.

Elsewhere, on 1 Corinthians 10:17, Chrysostom joins the one Eucharistic bread,
the one ecclesial body, and the very body adored by the Magi while it lay in the
manger. That is a genuine manger--Eucharist connection, although its immediate
argument is ethical and ecclesial: those nourished by the same body must live
as one.

\subsection{Augustine: believe, eat, abide}

Augustine's Tractates 25--27 on John resist two opposite reductions. His
celebrated command to believe and so eat does not make the sacrament
dispensable. It rejects the Capernaite imagination of cutting ordinary flesh
from a corpse. In the same tractates Augustine identifies the sacrament of
unity, interprets eating as abiding in Christ and Christ in the eater, and
warns that sacramental reception without inward participation does not yield
the promised life.

For Augustine, the sign and the reality belong together but are not
mechanically identical. Judas can receive the sacramental sign to judgment
without abiding in Christ. The life promised in John 6 belongs to fruitful
reception in faith and charity. Augustine therefore offers neither a merely
figurative denial nor an automatic sacramentalism. His ecclesial reading
presses the communicant toward incorporation into the Body signified and
received.

\subsection{Cyril of Alexandria: life-giving flesh}

Cyril of Alexandria reads John 6 through the Incarnation. The flesh is
life-giving not as an isolated material substance but because it is the flesh
of the Word, personally united to the giver of life. Eating and drinking join
the believer to Christ; the Eucharistic gift communicates life because the
one received is the incarnate Son. This logic guards both sides of the
chapter: divinity does not bypass flesh, and flesh is not treated apart from
the Word.

The available English Pusey--Randell witness is a nineteenth-century
translation and was not collated here against Pusey's Greek edition. This
study therefore paraphrases Cyril and does not make a wording-critical claim
from that translation.

\subsection{East and West: convergence without uniformity}

The representative Eastern witnesses emphasize distinct things. Ignatius
binds Eucharist to the flesh that suffered and rose; Irenaeus binds creation,
Eucharist, and bodily resurrection; the Jerusalem catechesis instructs
initiates; Chrysostom coordinates faith, mystery, union, and moral
transformation; Cyril of Alexandria explains life-giving flesh from the
hypostatic union of flesh and Word. They do not repeat one formula.

The representative Latin witnesses are likewise diverse. Cyprian places John
6 within daily bread, communion, and discipline. Ambrose contrasts manna's
mortality with the living Bread received as Christ's Body. Augustine stresses
faith, sacrament, ecclesial unity, abiding, and fruitful reception. Their
diversity is evidence of a living reception, not an embarrassment to be
flattened. The common center is the incarnate Christ as life-giving food; the
conceptual routes and pastoral pressures differ.

